\documentclass[a4paper,12pt]{article}
% packages!
\usepackage[utf8]{inputenc}
\usepackage{amsmath}
\usepackage{amssymb}% http://ctan.org/pkg/amssymb
\usepackage{booktabs}
\usepackage{color}
\usepackage{dirtree}
\usepackage{fancyhdr}
%\usepackage[cm]{fullpage}
\usepackage[width=14.5cm,
            height=25cm,
            top=2.5cm,
            inner=1.5cm,
            outer=1.5cm]{geometry}
\usepackage{graphicx}
\usepackage{grffile} % avoids errors when there are multiple dots in file names
\usepackage{listings}
\usepackage{longtable}
\usepackage{lscape}
\usepackage{lmodern}
\usepackage{makecell}
\usepackage{natbib}
\usepackage{pifont}% http://ctan.org/pkg/pifont
\usepackage{pbox}
\usepackage{sansmath}
\usepackage{setspace}
\usepackage{tabularx}
\usepackage{xcolor}
\usepackage{xparse}
% Alter some LaTeX defaults for better treatment of figures:
    \renewcommand{\topfraction}{0.9}    % max fraction of floats at top
    \renewcommand{\bottomfraction}{0.8} % max fraction of floats at bottom
    \setcounter{topnumber}{2}
    \setcounter{bottomnumber}{2}
    \setcounter{totalnumber}{4}     % 2 may work better
    \setcounter{dbltopnumber}{2}    % for 2-column pages
    \renewcommand{\dbltopfraction}{0.9} % fit big float above 2-col. text
    \renewcommand{\textfraction}{0.07}  % allow minimal text w. figs
    \renewcommand{\floatpagefraction}{0.7}      % require fuller float pages
    \renewcommand{\dblfloatpagefraction}{0.7}   % require fuller float pages

\renewcommand*\familydefault{\sfdefault}
\sansmath

\definecolor{darkgreen}{rgb}{0,0.6,0}
\definecolor{darkred}{rgb}{0.85,0,0}
\definecolor{gray}{rgb}{0.5,0.5,0.5}
\definecolor{mauve}{rgb}{0.58,0,0.82}
\definecolor{graybg}{rgb}{0.95,0.95,0.95}

\newcommand{\cmark}{\textcolor{darkgreen}{\ding{51}}}%
\newcommand{\xmark}{\textcolor{darkred}{\ding{55}}}%
\newcommand{\smark}{\ding{72}}%
\NewDocumentCommand{\codeword}{v}{%
\texttt{\textcolor{blue}{#1}}%
}

% set presets for Python code:
\lstdefinestyle{python}{
  language=Python
}
\lstdefinestyle{bash}{
  language=bash
}
\lstset{
   basicstyle={\footnotesize\ttfamily},
   showstringspaces=false,
   keywordstyle={\bfseries\color{blue}},
   numbers=none,
   numberstyle=\tiny\color{gray},
   commentstyle=\color{darkgreen},
   stringstyle=\color{mauve},
   breaklines=true,
   backgroundcolor=\color{graybg}
}

% Hurenkinder und Schusterjungen verhindern
\clubpenalty10000
\widowpenalty10000
\displaywidowpenalty=10000

%opening
\title{Documentation of the adapted Fractions Skill Score:\\
       Summed Area Table implementation (\texttt{fss\_SAT.py})}
\author{Phillip Scheffknecht}

\pagestyle{fancy}
\fancyhf{}
\rhead{\rightmark}
\lhead{FSS --- SAT method}
\rfoot{Page \thepage}

\begin{document}
\begin{spacing}{1.05}
\maketitle

\begin{abstract}
This document describes \texttt{fss\_SAT.py}, the summed area table (SAT)
implementation of the Fractions Skill Score (FSS). The implementation is based
on the reference code published by \citet{faggian2015}. We first recap the SAT
algorithm and its reference implementation, and then document the adaptations
and extensions made on top of it: a vectorised neighbourhood filter, a
BLAS-based score evaluation, additional threshold modes, percentile thresholds,
an over-estimation diagnostic, an ensemble (eFSS) variant, parallelisation over
thresholds, transparent handling of missing (NaN) data through per-window
valid-point weighting, and the continuous/weighted FSS (\texttt{CWFSS}). A final
section documents the automated cross-method test suite that guards these
adaptations.\\
\end{abstract}

\tableofcontents
\newpage

\section{Introduction}

The Fractions Skill Score (FSS) is a neighbourhood-based verification metric
for spatial forecasts, introduced by \citet{roberts2008}. Rather than comparing
forecast and observation values point by point, it compares the \emph{fraction}
of grid points that exceed a given threshold within a sliding window
(neighbourhood), as a function of both the threshold and the window
(neighbourhood) size. This avoids the \emph{double penalty} that afflicts
point-wise scores when a feature is forecast with a small spatial displacement
and reveals the spatial scale at which a forecast attains useful skill. It does
so by treating the spatial distribution of events within a window
probabilistically: a perfect forecast is one with the same event frequency as
the observation within the window, regardless of the exact placement.

For a single threshold and window size the score is \citep{roberts2008}
%
\begin{equation}
  \mathrm{FSS} = 1 -
  \frac{\dfrac{1}{N}\displaystyle\sum_{i=1}^{N}\bigl(p_f - p_o\bigr)^2}
       {\dfrac{1}{N}\displaystyle\sum_{i=1}^{N}p_f^{\,2}
        + \dfrac{1}{N}\displaystyle\sum_{i=1}^{N}p_o^{\,2}},
  \label{eq:fss}
\end{equation}
%
where $N$ is the number of windows in the domain (one centred on each grid
point), $p_f$ is the forecast fraction and $p_o$ the observed fraction of the
window. $\mathrm{FSS}=1$ denotes a perfect forecast and $\mathrm{FSS}=0$ no
skill.

The cost of the FSS comes from evaluating the windowed fractions $p_f$, $p_o$
for many window sizes and thresholds. \citet{faggian2015} showed that this
windowed summation can be computed very efficiently using a \emph{summed area
table} (also called an integral image): the table is built once per threshold
with two cumulative sums, after which the sum over any rectangular window is
obtained from just four array lookups, independent of the window size. This is
the algorithm implemented in \texttt{fss\_SAT.py}.

\subsection{Scope of this document}

\texttt{fss\_SAT.py} started from the reference Python implementation in the
appendix of \citet{faggian2015}. The purpose of this document is to record the
\emph{adaptations and extensions} that were made on top of that baseline. It is
organised as follows:

\begin{itemize}
  \item Section~\ref{sec:baseline} recaps the SAT algorithm of
        \citet{faggian2015} and its reference implementation, establishing the
        notation and the starting point for the changes.
  \item Section~\ref{sec:adaptations} documents each adaptation made in
        \texttt{fss\_SAT.py}: the performance rewrites that preserve the
        original result bit-for-bit, the feature extensions (threshold modes,
        percentile thresholds, ensemble FSS, parallelisation), and the
        substantive algorithmic additions (missing-data handling and the
        continuous/weighted FSS).
\end{itemize}

A companion FFT-based implementation (\texttt{fss\_FFT.py}) and an OpenCL
implementation (\texttt{fss\_OCL.py}) exist in the same repository and compute
the same score; they are out of scope here and documented separately.


\section{The summed area table baseline}
\label{sec:baseline}

This section recaps the algorithm of \citet{faggian2015} as far as it is needed
to follow the adaptations in Section~\ref{sec:adaptations}. The notation follows
the paper.

\subsection{Summed area table and windowed sums}

A field is first reduced to a binary event field $I$ by thresholding (for the
classic FSS, $I(i,j)=1$ where the field exceeds the threshold and $0$
otherwise). The summed area table $\hat{O}$ stores at each point the sum of all
cells above and to the left, inclusive:
%
\begin{equation}
  \hat{O}(i,j) = \sum_{\hat{\imath}\le i}\;\sum_{\hat{\jmath}\le j}
                 O(\hat{\imath},\hat{\jmath}).
  \label{eq:sat}
\end{equation}
%
In the code this is a double cumulative sum (\codeword{compute_integral_table}):

{\setstretch{1}\begin{lstlisting}[style=python]
def compute_integral_table(field):
    if field.ndim == 2:
        return field.cumsum(1).cumsum(0)
    elif field.ndim == 3:
        return field.cumsum(2).cumsum(1)
\end{lstlisting}}

Once $\hat{O}$ is available, the sum over a window of half-width $w$ centred on
$(i,j)$ is obtained from just four lookups,
%
\begin{equation}
  \sum_{\hat{\imath}=i-w}^{i+w}\;\sum_{\hat{\jmath}=j-w}^{j+w}
  O(\hat{\imath},\hat{\jmath})
  = \hat{O}(i+w, j+w) + \hat{O}(i-w, j-w)
  - \hat{O}(i-w, j+w) - \hat{O}(i+w, j-w),
  \label{eq:boxsum}
\end{equation}
%
\emph{independent of the window size}. This is the key property that makes the
SAT method cheap when many window sizes are evaluated per threshold: the table
is built once, and every window size is then four array lookups per grid point.
In the code the windowed sum is \codeword{integral_filter}.

\subsection{Boundary conditions}

Windows centred near the edge of the domain extend beyond the indexable region
of the table. \citet{faggian2015} clip the corner coordinates of
Eq.~\eqref{eq:boxsum} to the valid index range, which is how the reference
implementation (and \texttt{fss\_SAT.py}) handle the boundary. The reference
code realises the clip with \codeword{np.mgrid} and \codeword{np.clip}:

{\setstretch{1}\begin{lstlisting}[style=python]
r, c = np.mgrid[0:rows, 0:cols]
r0, c0 = (np.clip(r - w, 0, rows - 1), np.clip(c - w, 0, cols - 1))
r1, c1 = (np.clip(r + w, 0, rows - 1), np.clip(c + w, 0, cols - 1))
integral_table  = field[..., r1, c1] + field[..., r0, c0]
integral_table -= field[..., r0, c1] + field[..., r1, c0]
\end{lstlisting}}

\subsection{From windowed sums to the score}

Dividing each windowed sum by the window area $(2w+1)^2$ gives the fractions
$p_f$ and $p_o$ of Eq.~\eqref{eq:fss}. The reference \codeword{fss} routine
forms the fractions explicitly and evaluates the numerator and denominator with
\codeword{np.nanmean}:

{\setstretch{1}\begin{lstlisting}[style=python]
inv_window_area = 1. / (2. * w + 1.) ** 2
fhat = fhat * inv_window_area
ohat = ohat * inv_window_area
num   = np.nanmean(np.power(fhat - ohat, 2))
denom = np.nanmean(np.power(fhat, 2) + np.power(ohat, 2))
ret = num, denom, 1. - num / denom
\end{lstlisting}}

The higher-level routines (\codeword{fss_threshold}, \codeword{fss_cumsum_*},
\codeword{CWFSS}) build the binary tables once per threshold and call
\codeword{fss}/\codeword{integral_filter} for every requested window. This is
the structure that Section~\ref{sec:adaptations} modifies.


\section{Adaptations in \texttt{fss\_SAT.py}}
\label{sec:adaptations}

The changes fall into three groups: performance rewrites that leave the result
of \citet{faggian2015} unchanged (Sections~\ref{sec:filter}--\ref{sec:score}),
feature extensions (Sections~\ref{sec:modes}--\ref{sec:parallel}), and
algorithmic additions that go beyond the original method
(Sections~\ref{sec:nan}--\ref{sec:cwfss}). Table~\ref{tab:overview} gives an
overview.

\begin{table}[ht!]
\centering
\renewcommand{\arraystretch}{1.25}
\begin{tabularx}{\textwidth}{l l X}
\toprule
\textbf{Where} & \textbf{Change} & \textbf{Effect} \\
\midrule
\texttt{integral\_filter}   & \texttt{np.pad} + slicing   & same result, faster, no fancy indexing \\
\texttt{\_fss\_score}        & BLAS dot products           & same result, no temporary arrays \\
\texttt{\_build\_binary\_sat} & threshold modes            & \texttt{over}/\texttt{under}/\texttt{between}/\texttt{tolerance} \\
\texttt{fss\_threshold}      & percentile thresholds       & threshold given as a percentile of the field \\
\texttt{fss\_threshold}      & over-estimation \texttt{ovest} & frequency-bias diagnostic returned with FSS \\
\texttt{fss\_threshold\_eps} & ensemble FSS                & 3D forecast stack, exceedance probability \\
\texttt{fss\_cumsum\_parallel} & \texttt{joblib}           & threshold sweep optionally parallel (\texttt{n\_jobs}) \\
\texttt{\_fss\_score\_masked} & per-window valid weighting & transparent missing-data (NaN) handling \\
\texttt{CWFSS}               & R2 sampling, bootstrap, NaN & continuous/weighted FSS \\
\bottomrule
\end{tabularx}
\caption{Overview of the adaptations made on top of the \citet{faggian2015}
reference implementation.}
\label{tab:overview}
\end{table}

\subsection{Vectorised neighbourhood filter}
\label{sec:filter}

The reference \codeword{integral_filter} (Section~\ref{sec:baseline}) builds two
full coordinate grids with \codeword{np.mgrid}, clips them, and gathers the four
corner tables with fancy indexing. The rewrite achieves the identical clamped
box-sum of Eq.~\eqref{eq:boxsum} with edge-replicated padding and plain
contiguous slicing, which avoids the large index arrays and the gather:

{\setstretch{1}\begin{lstlisting}[style=python]
w = n // 2
if w < 1:
    return field
rows, cols = field.shape[-2], field.shape[-1]
D = 2 * w
p = np.pad(field, w, mode='edge')          # 2D case
return (p[D:D+rows, D:D+cols] + p[:rows, :cols]
      - p[:rows, D:D+cols] - p[D:D+rows, :cols])
\end{lstlisting}}

Padding the integral table with \codeword{mode='edge'} replicates the boundary
rows and columns, so a corner lookup that would have been clamped to the last
valid index instead reads the replicated edge value --- exactly what the
original \codeword{np.clip} produced. The boundary behaviour is therefore
\emph{identical} to the reference implementation; only the data access pattern
changes. The 3D (ensemble) case pads only the two trailing spatial axes.

\subsection{BLAS dot-product score}
\label{sec:score}

The reference \codeword{fss} scales both windowed fields to fractions and then
forms two temporary arrays, \codeword{(fhat-ohat)**2} and
\codeword{fhat**2+ohat**2}, before averaging. The rewrite (\codeword{_fss_score})
expresses the same quantities through three inner products of the
\emph{unscaled} box sums $S_f$, $S_o$, which BLAS evaluates without allocating
temporaries:
%
\begin{equation}
  \begin{aligned}
  \mathit{ff} &= \textstyle\sum_i S_f^2, &\quad
  \mathit{oo} &= \textstyle\sum_i S_o^2, &\quad
  \mathit{fo} &= \textstyle\sum_i S_f S_o .
  \end{aligned}
\end{equation}
%
Using the identity $\sum_i (S_f-S_o)^2 = \mathit{ff} + \mathit{oo} -
2\,\mathit{fo}$ and writing $A=(2w+1)^2$ for the window area, the numerator and
denominator of Eq.~\eqref{eq:fss} become
%
\begin{equation}
  \mathrm{num} = \frac{1}{A^2}\,\frac{\mathit{ff}+\mathit{oo}-2\,\mathit{fo}}{N},
  \qquad
  \mathrm{denom} = \frac{1}{A^2}\,\frac{\mathit{ff}+\mathit{oo}}{N},
\end{equation}
%
where $A^{-2}$ is the constant \codeword{inv_area_sq} factor (the fractions
$p=S/A$ enter squared). This is algebraically identical to the reference
result, but the score reduces to three \codeword{np.dot} calls plus scalar
arithmetic:

{\setstretch{1}\begin{lstlisting}[style=python]
ff = np.dot(fflat, fflat); oo = np.dot(oflat, oflat); fo = np.dot(fflat, oflat)
num   = inv_area_sq * (ff + oo - 2.0 * fo) / n
denom = inv_area_sq * (ff + oo) / n
if denom == 0.0:
    return 0.0, 0.0, np.nan
return num, denom, 1.0 - num / denom
\end{lstlisting}}

Note that \codeword{np.dot} is not NaN-aware, unlike the reference
\codeword{np.nanmean}. This fast path is therefore only taken when the fields
contain no missing data; the missing-data case is handled separately
(Section~\ref{sec:nan}).

\subsection{Threshold modes and percentile thresholds}
\label{sec:modes}

The reference implementation binarises with a single ``over'' comparison
($\text{field} > t$). \texttt{fss\_SAT.py} factors the binarisation into
\codeword{_build_binary_sat} and supports four \codeword{threshold_mode} values:
%
\begin{description}
  \item[\texttt{over}] $\text{field} > t$ (default, the classic FSS);
  \item[\texttt{under}] $\text{field} \le t$;
  \item[\texttt{between}] $t_1 < \text{field} \le t_2$, for verification of a
        value band (in the sweep, an extra lower edge $-1$ is inserted so that
        the lowest band still captures the zero values);
  \item[\texttt{tolerance}] $(1-\tau)\,t < \text{field} \le (1+\tau)\,t$,
        a relative tolerance band of half-width $\tau$ (\codeword{tolerance})
        around $t$.
\end{description}
%
Independently, \codeword{percentiles=True} interprets the threshold as a
percentile of the field rather than an absolute value; the forecast and
observation thresholds are then derived separately with \codeword{np.percentile}
(\codeword{np.nanpercentile} on the missing-data path).

\subsection{Over-estimation diagnostic}
\label{sec:ovest}

Alongside the numerator, denominator and score, \codeword{fss_threshold} now
returns an over-estimation (frequency-bias) value
%
\begin{equation}
  \mathrm{ovest} = \frac{\#\{\text{fcst} > t\} - \#\{\text{obs} > t\}}
                        {N_\text{points}},
\end{equation}
%
the difference between the forecast and observed event counts normalised by the
number of points. It is a domain-wide bias indicator and is broadcast to the
shape of the window axis so it travels through the same return structure as the
score. On the missing-data path it is computed over valid points only.

\subsection{Ensemble FSS (eFSS)}
\label{sec:eps}

\codeword{fss_threshold_eps} implements the ensemble extension of the FSS in the
spirit of \citet{duc2013}. The forecast is a 3D stack (members $\times$ rows
$\times$ cols). Instead of a binary forecast field, the per-point
\emph{exceedance probability} across members is used,
%
\begin{equation}
  p_f(i,j) = \frac{1}{M}\sum_{m=1}^{M}\mathbb{1}\!\left[\,\text{fcst}_m(i,j)
             \text{ exceeds } t\,\right],
\end{equation}
%
($M$ = number of members), and this probability field is fed into the same SAT
machinery as the deterministic case. The observation remains a single 2D field.
An assertion enforces the 3D forecast shape.

\subsection{Threshold sweep and parallelisation}
\label{sec:parallel}

\codeword{fss_cumsum_parallel} sweeps a list of thresholds, calling
\codeword{fss_threshold} (or \codeword{fss_threshold_eps} when
\codeword{eps=True}) once per threshold and reusing the cached integral tables
across all windows. The sweep can run serially (\codeword{n_jobs=1}) or in
parallel over thresholds through \codeword{joblib} (\codeword{n_jobs>1}):

{\setstretch{1}\begin{lstlisting}[style=python]
if n_jobs == 1:
    ret = [use_fss_threshold_func(fcst, obs, t1, t2, windows, ...) for t1, t2 in calls]
else:
    from joblib import Parallel, delayed
    ret = Parallel(n_jobs=n_jobs)(
        delayed(use_fss_threshold_func)(fcst, obs, t1, t2, windows, ...) for t1, t2 in calls)
\end{lstlisting}}

\codeword{fss_cumsum_frame} wraps the sweep and returns pandas
\codeword{DataFrame}s of numerator, denominator, score and over-estimation
(indexed by threshold, columns by window), or just the raw score array when
\codeword{raw=True}. Note that, because BLAS itself is multi-threaded,
oversubscription can occur when \codeword{n_jobs>1}; the benchmark
\codeword{bench_njobs.py} measures this with and without pinning the BLAS thread
count.

\subsection{Missing-data handling via per-window valid-point weighting}
\label{sec:nan}

This is the most substantial addition. The reference method has no concept of
missing data; a single NaN would propagate through the cumulative sums and
poison the whole table. \texttt{fss\_SAT.py} adds a mask-aware path that treats
a grid point as \emph{missing} when either the forecast \emph{or} the
observation is NaN there (\codeword{_validity_mask}). The clean path is taken
whenever the mask is all-valid, so fields without NaNs are unaffected and pay no
extra cost.

When NaNs are present, the binary fields are zeroed at missing points before the
SAT is built (\codeword{_build_binary_sat_masked}), so missing points never
contribute to a window sum. Each window is then weighted by its number of
\emph{valid} points,
%
\begin{equation}
  C = A - (\text{missing points in the window}),
  \label{eq:validcount}
\end{equation}
%
where $A=(2w+1)^2$ is the full window area and the missing-point count itself is
obtained from a SAT of the invalid mask (\codeword{_invalid_sat}) sampled with
the same \codeword{integral_filter}. With $S_f$, $S_o$ the box sums of the
mask-zeroed binary fields, the windowed fractions are $p_f=S_f/C$, $p_o=S_o/C$,
and weighting each window centre by $C$ collapses the weighted means to
%
\begin{equation}
  \mathrm{num} = \frac{\sum_w (S_f-S_o)^2 / C}{\sum_w C},
  \qquad
  \mathrm{denom} = \frac{\sum_w (S_f^2+S_o^2) / C}{\sum_w C}
  \label{eq:maskedscore}
\end{equation}
%
(\codeword{_fss_score_masked}). Windows with no valid points ($C=0$) drop out
naturally. The construction is deliberately consistent with the clean path: with
no missing data $C\equiv A$ everywhere, and Eq.~\eqref{eq:maskedscore} reduces
\emph{bit-for-bit} to the standard score of Section~\ref{sec:score}. Percentile
thresholds on this path use \codeword{np.nanpercentile}.

The same weighting is wired into the ensemble path (\codeword{fss_threshold_eps},
Section~\ref{sec:eps}): there a grid point counts when the observation is valid
\emph{and} at least one member is valid, and the exceedance probability is
averaged over the valid members only (\codeword{_eps_prob_masked}).

\subsection{Continuous/weighted FSS (\texttt{CWFSS})}
\label{sec:cwfss}

\codeword{CWFSS} aggregates the FSS over a \emph{range} of thresholds and
windows into a single scalar, instead of reporting the full
threshold\,$\times$\,window table. It draws \codeword{nsamples} (threshold,
window) pairs from the low-discrepancy R2 sequence (\codeword{R2}), which gives
even coverage of the 2D parameter space without a fixed grid, and computes the
FSS for each pair. Each sample is weighted by a threshold factor and a window
factor,
%
\begin{equation}
  \text{t\_factor} = \frac{t_\text{max}+t}{t_\text{max}},
  \qquad
  \text{w\_factor} = \frac{2\,w_\text{max}}{w_\text{max}+w},
\end{equation}
%
emphasising higher thresholds and smaller windows, and the weighted scores are
normalised by their achievable maximum to yield \codeword{self.cwfss}. A
\codeword{bootstrap} method resamples the stored per-sample scores to obtain a
confidence distribution.

\codeword{CWFSS} carries the same extensions as the batch functions. It honours
\codeword{threshold_mode}/\codeword{tolerance} through a shared
\codeword{_binarise} helper, and it is NaN-aware: threshold limits use
\codeword{np.nanmax}/\codeword{np.nanpercentile}, and when the fields contain
NaNs it uses the masked score of Eq.~\eqref{eq:maskedscore}. Samples whose FSS
is undefined (NaN, e.g.\ a window with no valid exceedance) are dropped from the
\codeword{np.nanmean}, and the theoretical maximum is masked identically so the
achievable ceiling stays at $1$ and is not lowered merely by the presence of
missing data.


\section{Verification and test suite}
\label{sec:tests}

The adaptations of Section~\ref{sec:adaptations} are guarded by an automated
\codeword{pytest} suite, \codeword{test_fss_consistency.py}. It serves two
purposes at once: it is a \emph{regression} suite that pins the behaviour of the
SAT code documented here, and a \emph{cross-method consistency} suite that checks
the SAT implementation against the two sibling backends in the repository --- the
FFT reference (\codeword{fss_FFT.py}) and the OpenCL implementation
(\codeword{fss_OCL.py}). Agreement between three independent code paths that
compute the same score is a strong correctness signal, and it is what lets the
performance rewrites of Sections~\ref{sec:filter}--\ref{sec:score} be trusted to
leave the result unchanged.

All tests run on a small set of reproducible Gaussian-bell fields built from a
fixed seed, so every value is deterministic. Two tolerances are used when
comparing methods, reflecting genuine algorithmic differences rather than bugs:

\begin{description}
  \item[$10^{-5}$] between SAT and OpenCL --- the \emph{same} algorithm, the only
        difference being \codeword{float32} arithmetic on the device versus
        \codeword{float64} on the host;
  \item[$10^{-3}$] between FFT and SAT --- the FFT convolution zero-pads at the
        domain boundary whereas the SAT filter edge-clamps
        (Section~\ref{sec:filter}), so the two agree in the interior but differ
        in a boundary frame.
\end{description}

The suite is run with

{\setstretch{1}\begin{lstlisting}[style=bash]
pytest test_fss_consistency.py -v
\end{lstlisting}}

\noindent and requires no arguments. The OpenCL tests detect whether a
\codeword{pyopencl} runtime is available and \emph{skip} themselves when it is
not, so the suite passes on a machine without a GPU/OpenCL stack.

\subsection{Coverage}
\label{sec:tests-coverage}

Table~\ref{tab:tests} groups the checks by theme. The right-hand column records
which of the three backends each theme exercises (S\,=\,SAT, F\,=\,FFT,
O\,=\,OpenCL).

\begin{table}[ht!]
\centering
\renewcommand{\arraystretch}{1.25}
\begin{tabularx}{\textwidth}{l X c}
\toprule
\textbf{Theme} & \textbf{What is verified} & \textbf{Backends} \\
\midrule
Cross-method
  & FFT, SAT and OpenCL scores agree within tolerance across thresholds
    $\times$ windows & S,F,O \\
Mathematical properties
  & perfect forecast $\to 1$, FSS $\in [0,1]$, NaN when no exceedances,
    monotonicity in window size, forecast/observation symmetry & S,F,O \\
Batch interfaces
  & \texttt{fss\_threshold}, \texttt{fss\_cumsum\_parallel} and
    \texttt{fss\_opencl\_arr} match the per-call single results & S,O \\
Regression guards
  & pinned reference scores for fixed (threshold, window) pairs, including
    every \texttt{threshold\_mode} & S \\
Threshold modes
  & \texttt{over}/\texttt{under}/\texttt{between}/\texttt{tolerance} preserve
    invariants and partition the domain correctly; the tolerance window is
    built from the forecast's own threshold on both bounds & S \\
\texttt{CWFSS}
  & R2 sequence, numerators / denominators / values / score and
    \texttt{bootstrap} are bit-identical to the reference, plus invariants
    and pinned values & S \\
Missing data
  & the masked path reduces \emph{exactly} to the clean path without NaNs, and
    preserves all invariants under injected NaNs (see
    Section~\ref{sec:tests-nan}) & S,F,O \\
\bottomrule
\end{tabularx}
\caption{Themes covered by \texttt{test\_fss\_consistency.py} and the backends
each one exercises (S\,=\,SAT, F\,=\,FFT, O\,=\,OpenCL).}
\label{tab:tests}
\end{table}

\subsection{The missing-data tests}
\label{sec:tests-nan}

Because the missing-data handling of Section~\ref{sec:nan} is the most
substantial addition, it carries the most checks, and they are mirrored across
all three backends (deterministic and ensemble for SAT and FFT; deterministic
for OpenCL, which has no ensemble entry point). The key guarantees are:

\begin{description}
  \item[Reduction to the clean path.] On NaN-free fields the masked score of
        Eq.~\eqref{eq:maskedscore} is compared against the ordinary fast path.
        For SAT and FFT this matches to machine precision ($\sim 10^{-9}$ or
        better), confirming that the valid count $C$ of
        Eq.~\eqref{eq:validcount} collapses to the constant window area $A$ when
        nothing is missing; for OpenCL it matches to the \codeword{float32}
        tolerance $10^{-5}$. A companion test verifies directly that
        $C \equiv A$ everywhere in the absence of NaNs.
  \item[Invariants under NaN.] With NaNs injected into both fields, the score
        stays in $[0,1]$, a perfect forecast (sharing the observation's NaNs)
        still scores $1$, the forecast/observation symmetry holds, and an
        all-missing field yields NaN rather than a spurious number.
  \item[Weighting semantics.] A hand-rolled weighted mean reproduces
        \codeword{_fss_score_masked}, pinning the per-window valid-point
        weighting of Eq.~\eqref{eq:maskedscore}.
  \item[Ensemble specifics.] A member that is entirely NaN does not change the
        score (the average is over valid members only), and the
        \codeword{nanpercentile} threshold path stays finite.
  \item[Cross-method agreement.] The FFT and OpenCL masked paths are checked
        against the SAT masked path. OpenCL also edge-clamps, so it agrees with
        SAT as closely as on clean data; FFT zero-pads, so it agrees within the
        interior at the usual $10^{-3}$ boundary tolerance.
  \item[Auto-dispatch.] Calling the public entry points
        (\codeword{fss_threshold}, \codeword{fourier_fss},
        \codeword{fourier_fss_eps}, \codeword{fss_opencl}) on a field that
        contains a NaN routes to the masked path automatically --- important for
        the FFT, where a raw NaN passed through \codeword{fftconvolve} would
        otherwise spread across the entire output.
\end{description}

\subsection{Result}
\label{sec:tests-result}

The suite currently comprises \textbf{308 tests}. With an OpenCL runtime present
all \textbf{308 pass}; on a machine without OpenCL the \textbf{60} device-backed
tests skip and the remaining \textbf{248 pass}. A representative run completes in
roughly six seconds.


\begin{thebibliography}{9}
\bibitem[Faggian et al.(2015)]{faggian2015}
Faggian, N., Roux, B., Steinle, P. and Ebert, B., 2015:
\textit{Fast calculation of the fractions skill score.}
MAUSAM, \textbf{66}, 3, 457--466.

\bibitem[Roberts and Lean(2008)]{roberts2008}
Roberts, N. M. and Lean, H. W., 2008:
\textit{Scale-selective verification of rainfall accumulations from
high-resolution forecasts of convective events.}
Mon. Wea. Rev., \textbf{136}, 78--97.

\bibitem[Duc et al.(2013)]{duc2013}
Duc, L., Saito, K. and Seko, H., 2013:
\textit{Spatial-temporal fractions verification for high-resolution ensemble
forecasts.} Tellus A, \textbf{65}.
\end{thebibliography}

\end{spacing}
\end{document}
